\documentclass[uplatex,dvipdfmx,11pt,a4paper]{bxjsarticle}
\usepackage{geometry}
\usepackage{amsmath,amssymb,amsfonts,bm}
\usepackage{booktabs}
\usepackage{longtable}
\usepackage{hyperref}
\usepackage{xcolor}
\usepackage{enumitem}
\geometry{margin=25mm}
\hypersetup{colorlinks=true,linkcolor=blue,citecolor=blue,urlcolor=blue}

\title{自社DSPにおける Contextual Bandit 広告絞り込みの問題定式化とアルゴリズム候補}
\author{Matching Algorithm Research Note}
\date{\today}

\newcommand{\R}{\mathbb{R}}
\newcommand{\E}{\mathbb{E}}
\newcommand{\ind}{\mathbf{1}}
\newcommand{\calA}{\mathcal{A}}
\newcommand{\calX}{\mathcal{X}}
\newcommand{\calD}{\mathcal{D}}

\begin{document}
\maketitle

\begin{abstract}
本稿では、決済後広告の配信問題を、自社DSP（Demand-Side Platform）における
Contextual Bandit による広告絞り込み問題として定式化する。ここでの目的は、
リアルタイム入札（RTB）をエンドツーエンドで解くことではなく、探索、
コールドスタート、遅延フィードバック、オークションに起因するセレクションバイアスを
考慮しながら、大量の広告候補から有望な少数の広告を選ぶ第一層を構築することである。
本稿では、意思決定プロセス、報酬、ログ要件、オフライン評価手順を定義し、
Disjoint LinUCB、Hybrid LinUCB、Linear Thompson Sampling、Logistic/GLM Bandit、
Neural Bandit などの Contextual Bandit アルゴリズムを比較する。最初に実装すべき
方針として、特徴量ベースの LinUCB / LinTS ベースライン、逆傾向スコアを意識した
ログ設計、そして hybrid 化および一般化線形モデルへの拡張経路を推奨する。
\end{abstract}

\section{対象範囲}
対象となるプロダクト面は、レンタカー予約サービスにおける決済後ページである。
ユーザーは予約を完了した直後であるため、システムは車種、地域、レンタル日数、
価格、デバイス、場合によってはセッション履歴などの文脈変数を観測できる。
DSP は広告集合を保有しており、各広告にはカテゴリ、対象地域、対象車種、広告主、
予算、コールドスタート状態などのメタデータが付与されている。絞り込み層は、
入札層へ渡す少数の広告集合を選択する。入札層は、その広告集合に対して価値予測、
予算ペーシング、オークションロジックを用いて入札価格を決定する。

本稿は絞り込み層に焦点を当てる。RTB には打ち切られたフィードバック、予算制約、
ペーシング、戦略的なオークション効果が含まれており、標準的な Contextual Bandit だけでは
完全には表現できないため、絞り込みと入札を意図的に分離して扱う。

\section{意思決定問題}
\subsection{対象オブジェクト}
インタラクション $t=1,2,\ldots,T$ において、以下を定義する。
\begin{itemize}[leftmargin=2em]
  \item $c_t \in \mathcal{C}$ は観測された購入文脈である。車種、地域、地域タイプ、
  レンタル日数、価格帯、デバイス、セッション単位の集約特徴、時刻特徴などを含む。
  \item $\mathcal{A}_t = \{a_{t1},\ldots,a_{tN_t}\}$ は、ビジネス上のハードフィルタを
  通過した広告集合である。ハードフィルタには、キャンペーンの有効状態、地域適合性、
  クリエイティブ審査、予算残高などが含まれる。
  \item $z(a) \in \mathcal{Z}$ は広告側特徴である。カテゴリ、対象地域、対象車種、
  広告主、キャンペーン経過日数、新規フラグ、利用可能であれば過去実績に基づく事前情報を含む。
  \item $x_t(a)=\phi(c_t,z(a))\in\R^d$ は、文脈と広告を結合した特徴ベクトルである。
  \item 絞り込みポリシーは $S_t\subseteq\mathcal{A}_t$ かつ $|S_t|=K$ となる部分集合を選び、
  入札層は $S_t$ を受け取る。
\end{itemize}

\subsection{報酬}
$Y_t(a)$ を、文脈 $c_t$ において広告 $a$ を表示した場合の潜在アウトカムとする。
プロダクト目標に応じて、報酬は例えば以下のように定義できる。
\begin{align}
  r_t(a) &= \ind\{\text{クリック}\}, \\
  r_t(a) &= v_{\mathrm{cv}}\ind\{\text{コンバージョン}\}, \\
  r_t(a) &= \alpha_{\mathrm{ctr}}\ind\{\text{クリック}\}
          + \alpha_{\mathrm{cvr}}v_{\mathrm{cv}}\ind\{\text{コンバージョン}\}
          - \alpha_{\mathrm{cost}}\,\mathrm{cost}.
\end{align}
最初の絞り込み層プロトタイプでは、クリック報酬を用いてよい。クリックは頻度が高く、
比較的早く観測できるためである。本番最適化では、遅延コンバージョンとオークション費用を
考慮したうえで、インプレッションあたり価値、または落札インプレッションあたり価値へ
報酬を移行すべきである。

Top-$K$ 絞り込みの理想目的は次のように書ける。
\begin{equation}
  \max_{\pi}\; \E\left[\sum_{t=1}^{T}\sum_{a\in S_t^\pi} r_t(a)\right],
  \quad S_t^\pi = \pi(c_t,\mathcal{A}_t),\quad |S_t^\pi|=K.
\end{equation}
同値な見方として、期待報酬を知っている oracle に対する regret は次のように定義できる。
\begin{equation}
  R_T = \sum_{t=1}^{T}
  \left(
    \sum_{a\in S_t^\star}\mu(c_t,a)
    - \sum_{a\in S_t^\pi}\mu(c_t,a)
  \right),
\end{equation}
ここで $\mu(c,a)=\E[r\mid c,a]$ であり、$S_t^\star$ は $\mu$ による上位 $K$ 件の広告集合である。

\subsection{特徴量設計}
線形モデルを最初に用いる場合、交互作用は明示的に特徴量へ入れる必要がある。
最小構成の特徴写像は次のようにできる。
\begin{align}
  x_t(a)=\big[&\mathrm{onehot}(\mathrm{car\_type}),
  \mathrm{onehot}(\mathrm{area}),
  \mathrm{onehot}(\mathrm{area\_type}),
  \mathrm{price\_bucket},
  \mathrm{rental\_days\_scaled},\\
  &\mathrm{onehot}(\mathrm{ad\_category}),
  \mathrm{is\_new},
  \ind\{\mathrm{context.area}=\mathrm{ad.target\_area}\},\\
  &\ind\{\mathrm{area\_type=resort}\}\times
    \ind\{\mathrm{category}\in\{\mathrm{spot},\mathrm{restaurant}\}\},\\
  &\ind\{\mathrm{car\_type}=\mathrm{ad.target\_car\_type}\}\times
    \ind\{\mathrm{category=goods}\},\\
  &\mathrm{car\_type}\times\mathrm{ad\_category},
  \mathrm{area}\times\mathrm{ad\_category}
  \big].
\end{align}
特徴写像はバージョン管理すべきである。ログに残す各意思決定には、特徴量バージョン、
候補広告集合、選択広告集合、ランキングスコア、行動確率または傾向スコアを保存する必要がある。

\section{フィードバックとバイアス}
RTB ではフィードバックが打ち切られる。システムがクリックやコンバージョンを観測できるのは、
広告が選択され、入札され、落札され、実際に表示された場合に限られる。したがって、
落札インプレッションだけで素朴に学習すると、広告そのものの関連度ではなく、
オークションやペーシングシステムの挙動を学習してしまう可能性がある。

本番品質のログには、以下を含めるべきである。
\begin{itemize}[leftmargin=2em]
  \item リクエストID、文脈、タイムスタンプ、プライバシーに配慮したユーザーまたはセッションキー。
  \item 絞り込み前の全候補広告。
  \item 選択された top-$K$ 集合、スコア、傾向スコア。
  \item 入札判断、入札価格、ペーシング係数、shading 係数。
  \item 勝敗、利用可能であれば清算価格、インプレッション、クリック、コンバージョン、遅延ラベル。
  \item モデルバージョン、特徴量バージョン、探索パラメータ。
\end{itemize}

オフライン評価には、ランダム化された探索トラフィックと、IPS（Inverse Propensity Scoring）、
自己正規化IPS、Doubly Robust 推定量などを用いる。貪欲方策から得られた決定論的ログだけでは、
代替ポリシーを不偏に評価するには不十分である。

\section{アルゴリズム候補}
\subsection{Disjoint LinUCB}
古典的な Disjoint LinUCB は、腕ごとに独立したパラメータを保持する。腕 $a$ について、
\begin{align}
  A_a &= \lambda I + \sum_{s:a_s=a} x_s x_s^\top, &
  b_a &= \sum_{s:a_s=a} r_s x_s, &
  \hat\theta_a &= A_a^{-1}b_a.
\end{align}
スコアは次で与えられる。
\begin{equation}
  p_t(a)=\hat\theta_a^\top x_t(a)+\alpha\sqrt{x_t(a)^\top A_a^{-1}x_t(a)}.
\end{equation}
実装が容易で解釈しやすい一方、完全な腕別学習は新規広告に弱い。新しく導入された腕には
履歴がないためである。

\subsection{特徴量ベースの共有 LinUCB}
実用的な第一歩は、広告を特徴量で表し、1つの共有モデルを学習することである。
\begin{align}
  A &= \lambda I + \sum_s x_s x_s^\top, &
  b &= \sum_s r_s x_s, &
  \hat\theta &= A^{-1}b,\\
  p_t(a)&=\hat\theta^\top x_t(a)+\alpha\sqrt{x_t(a)^\top A^{-1}x_t(a)}.
\end{align}
この方法では、メタデータと交互作用特徴を通じて広告間で学習を転移できるため、
コールドスタートに強い。本リポジトリで最初に信頼できるベースラインとして実装すべきである。

\subsection{Hybrid LinUCB}
Hybrid LinUCB は、報酬を全体共通成分と腕固有成分に分解する。
\begin{equation}
  \E[r_t(a)] \approx z_t(a)^\top\beta + x_t(a)^\top\theta_a.
\end{equation}
ここで $z_t(a)$ は腕をまたいで共有される特徴を表し、$x_t(a)$ は腕固有の特徴を表す。
広告主や広告単位の固有効果を学習するだけの十分なデータがあり、かつコールドスタートへの
一般化を失いたくない段階で適している。

\subsection{Linear Thompson Sampling}
Linear Thompson Sampling は、LinUCB と同じベイズ線形回帰構造を用いるが、楽観的ボーナスを
足す代わりにパラメータベクトルをサンプリングする。
\begin{align}
  \hat\theta &= A^{-1}b, &
  \tilde\theta_t &\sim \mathcal{N}(\hat\theta, v^2 A^{-1}), &
  p_t(a)&=\tilde\theta_t^\top x_t(a).
\end{align}
本番ランキングシステムでは、UCB よりも滑らかに振る舞うことが多い。線形代数部分の多くを
LinUCB と共有できるため、主要な比較ベースラインとして実装すべきである。

\subsection{Logistic または GLM Bandit}
クリックやコンバージョンは二値であるため、線形平均よりも logistic 平均関数のほうが自然である。
\begin{equation}
  \Pr(r=1\mid x,a)=\sigma(\theta^\top x), \quad
  \sigma(u)=\frac{1}{1+e^{-u}}.
\end{equation}
特に低CTR領域で線形予測の較正が悪い場合、GLM Bandit が有効である。ただし、オンラインの
信頼区間には曲率や近似事後更新が必要になるため、実装はより複雑になる。

\subsection{Neural Bandit}
Neural Bandit は、手作業で設計した線形特徴写像をニューラル表現 $f_\omega(c,a)$ に置き換える。
非線形な交互作用を捉えられる一方で、より多くのデータ、慎重な不確実性推定、強い監視が必要である。
最初の本番アルゴリズムとしては推奨しない。

\section{比較表}
\begin{longtable}{p{0.20\linewidth}p{0.24\linewidth}p{0.24\linewidth}p{0.20\linewidth}}
\toprule
手法 & 強み & 弱み & 推奨される役割 \\
\midrule
Disjoint LinUCB & 単純で解釈しやすい。広告ごとのパラメータを持てる。 & コールドスタートに弱い。広告数が多いと多数の行列が必要。 & 教育用ベースライン \\
共有 LinUCB & 単純、高速、広告間で転移可能、コールドスタートに対応しやすい。 & 線形仮定がある。特徴設計が必要。 & 最初の実装 \\
Hybrid LinUCB & 全体転移と腕固有効果のバランスを取れる。 & 管理対象とパラメータが増える。 & 第2段階の拡張 \\
Linear Thompson Sampling & 滑らかな探索が可能。LinUCB と状態を共有しやすい。 & サンプリングと分散較正が必要。 & LinUCB との主要A/B比較 \\
Logistic/GLM Bandit & 二値報酬と較正に適している。 & オンライン推論がより複雑。 & 線形ベースライン診断後に導入 \\
Neural Bandit & 非線形表現力が高い。 & データ要求が大きく、不確実性推定が難しい。 & 後続の研究テーマ \\
\bottomrule
\end{longtable}

\section{推奨実装計画}
\begin{enumerate}[leftmargin=2em]
  \item 文脈、広告、候補集合、行動、傾向スコア、アウトカムに関する本番品質のスキーマを定義する。
  \item $A^{-1}$ の Sherman--Morrison 更新を持つ共有線形モデル抽象を実装する。
  \item 同一の特徴量ストア上に LinUCB と Linear Thompson Sampling の両方を実装する。
  \item Semi-bandit 型の top-$K$ 学習を追加し、表示された各広告について観測報酬でモデルを更新する。
  \item オフラインで効果を主張する前に、ランダム化ログトラフィックと傾向スコア記録を追加する。
  \item シミュレーションでは reward@K、CTR@K、累積 regret を評価し、ログ上では IPS/SNIPS と較正プロットを用いる。
  \item 共有モデルで安定したログと十分な広告別観測数が得られてから Hybrid LinUCB を導入する。
  \item 入札層は分離しておく。絞り込み層からの期待価値は、ペーシング、予算、bid shading とともに
  入札価格を決める入力の1つにする。
\end{enumerate}

\section{参考文献}
\begin{thebibliography}{9}
\bibitem{li2010linucb}
L. Li, W. Chu, J. Langford, and R. E. Schapire.
\newblock A Contextual-Bandit Approach to Personalized News Article Recommendation.
\newblock WWW, 2010.

\bibitem{chu2011linear}
W. Chu, L. Li, L. Reyzin, and R. E. Schapire.
\newblock Contextual Bandits with Linear Payoff Functions.
\newblock AISTATS, 2011.

\bibitem{agrawal2013ts}
S. Agrawal and N. Goyal.
\newblock Thompson Sampling for Contextual Bandits with Linear Payoffs.
\newblock ICML, 2013.

\bibitem{li2011offline}
L. Li, W. Chu, J. Langford, and X. Wang.
\newblock Unbiased Offline Evaluation of Contextual-Bandit-based News Article Recommendation Algorithms.
\newblock WSDM, 2011.

\bibitem{wang2017rtb}
J. Wang, W. Zhang, and S. Yuan.
\newblock Display Advertising with Real-Time Bidding (RTB) and Behavioural Targeting.
\newblock Foundations and Trends in Information Retrieval, 2017.

\bibitem{zhang2014optimal}
W. Zhang, S. Yuan, and J. Wang.
\newblock Optimal Real-Time Bidding for Display Advertising.
\newblock KDD, 2014.

\bibitem{xu2015pacing}
J. Xu, K.-c. Lee, W. Li, Q. Qi, and Q. Lu.
\newblock Smart Pacing for Effective Online Ad Campaign Optimization.
\newblock KDD, 2015.

\bibitem{zhang2016censored}
W. Zhang, T. Zhou, J. Wang, and J. Xu.
\newblock Bid-aware Gradient Descent for Unbiased Learning with Censored Data in Display Advertising.
\newblock KDD, 2016.
\end{thebibliography}

\end{document}
